\documentclass[11pt,reqno]{amsart}

\usepackage{amsmath,amssymb,amsthm,mathtools}
\usepackage{tikz}
\usetikzlibrary{shapes.geometric, positioning, arrows.meta, decorations.pathmorphing, patterns}
\usepackage{array,booktabs,tabularx,float}
\usepackage[hidelinks]{hyperref}

\allowdisplaybreaks
\raggedbottom
\setlength{\emergencystretch}{2em}

\newtheorem{theorem}{Theorem}[section]
\newtheorem{proposition}[theorem]{Proposition}
\newtheorem{lemma}[theorem]{Lemma}
\newtheorem{corollary}[theorem]{Corollary}
\theoremstyle{definition}
\newtheorem{definition}[theorem]{Definition}
\newtheorem{example}[theorem]{Example}
\newtheorem{warning}[theorem]{Warning}
\newtheorem{openproblem}[theorem]{Open Problem}
\theoremstyle{remark}
\newtheorem{remark}[theorem]{Remark}

\newcommand{\C}{\mathbb C}
\newcommand{\Z}{\mathbb Z}
\newcommand{\Chat}{\widehat{\mathbb C}}
\newcommand{\Fill}{\operatorname{Fill}}
\newcommand{\Exp}{\operatorname{Exp}}
\newcommand{\Int}{\operatorname{int}}
\newcommand{\dist}{\operatorname{dist}}
\newcommand{\id}{\operatorname{id}}
\newcommand{\Om}{\Omega}
\newcommand{\Omfree}{\Omega_{\mathrm{free}}}
\newcolumntype{P}[1]{>{\raggedright\arraybackslash}p{#1}}
\newcolumntype{Y}{>{\raggedright\arraybackslash}X}

\title[Protected planar carriers]{Protected planar carriers:\\
separation ledgers and access-preserving routing}
\author{Matthew Robert Wesney}
\date{Draft II of 2026-07-27}
\subjclass[2020]{Primary 54F15, 30C20; Secondary 37F20, 55M05}
\keywords{plane continuum, prime end, Alexander duality, protected access,
connector, filled hull, Hausdorff limit}
\hypersetup{
  pdftitle={Protected planar carriers: separation ledgers and access-preserving routing},
  pdfauthor={Matthew Robert Wesney},
  pdfsubject={Planar continuum theory and access-preserving routing},
  pdfkeywords={plane continuum, prime end, Alexander duality, protected access,
    connector, filled hull, Hausdorff limit}
}

\begin{document}

\begin{abstract}
We give a finite planar theory for joining pairwise disjoint full continua
while preserving specified prime-end accesses and avoiding a frozen compact
inventory.  A protected ribbon-routing certificate is realized by a full
carrier whose null crosscut chains survive, remain distinct when their
relative collar tails are disjoint, and retain cyclic order, without any
local-connectivity hypothesis.  Certificate existence is equivalent to one
component of the frozen free domain admitting simultaneous singleton ports
to all atoms.  Packet and cellwise criteria give checkable sufficient
conditions, while fixed-assignment audits give persistent witnesses for
merger, escape, and port loss.

We also prove an attachment exact sequence.  If a connected connector $C$
meets a full atom $A$ in $T$, then the new first \v{C}ech cohomology fits into
the short exact sequence
$0\to\widetilde{\check H}^{\,0}(T)\to
\check H^1(C\cup A)\to
\ker[\check H^1(C)\to\check H^1(T)]\to0$.
Iteration gives the separation ledger
\[
\sigma\!\left(C\cup\bigcup_iA_i\right)
=\sigma(C)+\sum_i
\bigl(\#\pi_0(C\cap A_i)-1-\sigma(C\cap A_i)\bigr)
\]
whenever the ranks are finite.  Thus disconnected traces create holes and
separating traces fill pre-existing holes.  For a full connector the union
is full exactly when every attachment trace is connected.  Arbitrarily many
holes can nevertheless be born in a Hausdorff limit of full tree carriers
with singleton traces.  By contrast, decreasing full connector envelopes
with nonempty connected traces are stable and preserve common protected
data.
These results isolate the precise finite/infinite compactness gate; they make
no claim of constructing a virtual dynamical carrier or proving Mandelbrot
local connectivity.
\end{abstract}

\maketitle
\tableofcontents

%%INSERT A:intro_v2

%%INSERT A:preliminaries

%%INSERT B:obstructions

%%INSERT A:comparison

%%INSERT C:certificates

%%INSERT G:separation_budget

%%INSERT D:criterion

%%INSERT E:packets

%%INSERT E:audit

%%INSERT G:scope_v2

\section*{Acknowledgments}
AI-assisted drafting and citation checking were used.  The author
independently verified and assumes responsibility for all mathematical
statements, proofs, and references.

\appendix

\section{Provenance and statement crosswalk}\label{sec:crosswalk}

The certificate, criterion, and packet results originated in the source
research manuscript \cite{MWman}.  This draft corrects the slit statement,
uses null crosscut chains and relative collar tails throughout, repairs the
single-atom construction, and adds the separation ledger, the sharp
Hausdorff-limit example, and nested-limit stability.  No statement labeled
open or conditional in the source has been promoted.

\begin{center}
\small
\begin{tabularx}{\textwidth}{@{}P{0.35\textwidth}P{0.22\textwidth}Y@{}}
\toprule
Present result & Source & Status in this paper\\
\midrule
Theorem~\ref{thm:no-least} & Theorem 6.12 & extracted\\
Proposition~\ref{prop:access-collision} & Proposition 6.15 & extracted\\
Definition~\ref{def:certificate} & Definition 10.1 &
relative-tail and incidence conventions made explicit\\
Lemma~\ref{lem:tree-full} & Lemma 10.2 & two proofs supplied\\
Lemma~\ref{lem:crosscut-survival} & Lemma 10.3 &
relative-tail formulation\\
Theorem~\ref{thm:realization} & Theorem 10.4 &
proof synchronized with corrected conventions\\
Proposition~\ref{prop:naturality} & Proposition 10.5 & extracted\\
Theorem~\ref{thm:separation-budget} & Draft II &
classical duality/Mayer--Vietoris consequence, packaged for carriers\\
Corollary~\ref{cor:full-connector-traces} & Draft II &
classical zero-defect consequence, packaged for carriers\\
Proposition~\ref{prop:trace-birth} & Draft II &
Hausdorff-limit separation example proved in this draft\\
Theorem~\ref{thm:nested-limit-stability} & Draft II &
decreasing-envelope theorem proved in this draft\\
Lemma~\ref{lem:slits} & Lemma 10.7 &
repaired; see Remark~\ref{rem:repair}\\
Theorem~\ref{thm:free-component} & Theorem 10.8 &
proof rebuilt around the repaired slit lemma; $s=1$ included\\
Definition~\ref{def:packet} & Definition 10.9 &
compact atom--tail blocks and named compact obstacle pieces\\
Lemma~\ref{lem:packets} & Lemma 10.10 &
compacta duality made explicit\\
Theorem~\ref{thm:common-domain} & Theorem 10.11 & extracted\\
Lemma~\ref{lem:fill-absorption} & Lemma 10.12 & extracted\\
Lemma~\ref{lem:port-localization} & Lemma 10.13 & extracted\\
Theorem~\ref{thm:audit} & Theorem 10.14 & extracted\\
Proposition~\ref{prop:cellwise} & Proposition 10.15 & extracted\\
\bottomrule
\end{tabularx}
\end{center}

Examples~\ref{ex:no-common-component} and~\ref{ex:merger}, together with
Section~\ref{sec:comparison}, are expository additions.  The
direct-decoration application formerly included in Draft~I has been removed
because its strengthened form is the subject of the companion paper
\cite{Wes26}.

\begin{thebibliography}{99}

\bibitem{Bor31}
K.~Borsuk,
\emph{Quelques th\'eor\`emes sur les ensembles unicoh\'erents},
Fund. Math. \textbf{17} (1931), 171--209.

\bibitem{DH85}
A.~Douady and J.~H.~Hubbard,
\emph{On the dynamics of polynomial-like mappings},
Ann. Sci. \'Ec. Norm. Sup\'er. (4) \textbf{18} (1985), 287--343,
doi:10.24033/asens.1491.

\bibitem{DHOrsay}
A.~Douady and J.~H.~Hubbard,
\emph{\'Etude dynamique des polyn\^omes complexes, Parties I--II},
Publications Math\'ematiques d'Orsay 84-02 (1984) and 85-04 (1985).

\bibitem{Dud25}
D.~Dudko,
\emph{On the MLC conjecture and the renormalization theory in complex
dynamics},
to appear in the Proceedings of the ICM 2026,
arXiv:2512.24171v1, submitted 30 December 2025.

\bibitem{Eil36}
S.~Eilenberg,
\emph{Sur les espaces multicoh\'erents. I},
Fund. Math. \textbf{27} (1936), 153--190.

\bibitem{Hat02}
A.~Hatcher,
\emph{Algebraic Topology},
Cambridge University Press, Cambridge, 2002.

\bibitem{KKL26}
J.~Kahn, A.~Kapiamba, and M.~Lyubich,
\emph{MLC for parabolically bounded primitive renormalization},
arXiv:2606.27272v1, submitted 25 June 2026.

\bibitem{KL08}
J.~Kahn and M.~Lyubich,
\emph{A priori bounds for some infinitely renormalizable quadratics.
II. Decorations},
Ann. Sci. \'Ec. Norm. Sup\'er. (4) \textbf{41} (2008), 57--84,
doi:10.24033/asens.2063.

\bibitem{Jan15}
Z.~Janiszewski,
\emph{O rozcinaniu p\l aszczyzny przez kontinua}
[French summary: \emph{Sur les coupures du plan faites par des continus}],
Prace Matematyczno-Fizyczne \textbf{26} (1915), 11--63.

\bibitem{Moo25}
R.~L.~Moore,
\emph{Concerning upper semi-continuous collections of continua},
Trans. Amer. Math. Soc. \textbf{27} (1925), 416--428,
doi:10.1090/S0002-9947-1925-1501320-8.

\bibitem{Mul22}
A.~M.~Mullikin,
\emph{Certain theorems relating to plane connected point sets},
Trans. Amer. Math. Soc. \textbf{24} (1922), 144--162,
doi:10.2307/1989040.

\bibitem{MWman}
M.~R.~Wesney,
\emph{A marked-width reduction of Mandelbrot local connectivity},
unpublished research manuscript, 2026 (build of 2026-07-27).

\bibitem{New51}
M.~H.~A.~Newman,
\emph{Elements of the Topology of Plane Sets of Points},
2nd ed., Cambridge University Press, Cambridge, 1951.

\bibitem{Pom92}
Ch.~Pommerenke,
\emph{Boundary Behaviour of Conformal Maps},
Grundlehren der mathematischen Wissenschaften \textbf{299},
Springer-Verlag, Berlin, 1992.

\bibitem{Spa66}
E.~H.~Spanier,
\emph{Algebraic Topology},
McGraw--Hill, New York, 1966.

\bibitem{Wes26}
M.~R.~Wesney,
\emph{The direct-decoration rooted-end compactification},
companion manuscript in preparation, 2026.

\bibitem{Why42}
G.~T.~Whyburn,
\emph{Analytic Topology},
American Mathematical Society Colloquium Publications \textbf{28},
American Mathematical Society, New York, 1942.

\end{thebibliography}

\end{document}
